Microservice architectures promise that each service can change and be deployed without depending on the others. When several services reuse the same internal component, however, a change to that component may force them to change together. Practitioners have reported this risk, but it has rarely been measured in real systems. This work analyzes how sharing components affects the evolutionary autonomy of microservices, using the code and history of software repositories. From two public datasets with 395 candidates, eight open-source Java systems were randomly selected and mined with PyDriller. Co-change between services, measured by the Jaccard index over commits, was crossed with the dependencies declared in build files, and the 41 shared components found were classified by content. About half of the examined systems with three or more services share code among them, and in the median system 87\% of the services depend on some shared component. In five of the six systems where the comparison was possible, pairs of services sharing a component changed together more often, with a significant difference and a large effect in two of them. Manual inspection of 152 sampled commits showed, however, that only 42\% of the co-changes reflect a real dependency, while most of the rest come from batch maintenance. The share of a component's changes that reached its dependent services was not related to the number of services using it, nor to whether its code is technical or domain-specific. The highest rates appeared in small modules that describe a service's interface. The results suggest that the risk of a shared component depends less on its reach than on its role as an agreement between services.

\keywords{Microservices. Software reuse. Logical coupling. Mining software repositories. Shared components}
